% Chapter Template

\section{Internal Validity} % Main chapter title

\label{Sensitivity Analysis} % Change X to a consecutive number; for referencing this chapter elsewhere, use \ref{ChapterX}
In this section we take the same approach as \cite{caiani_agent_2016} to analyze the internal dynamics of the optimal consumption agent based model. We will compute the auto-correlations and the cross-correlations of the variables that are generated by the model and compare these with their empirical counterparts\footnote{The model is not built for long-term forecasts. For instance, the assumption that debt is repaid with a constant repayment rate results in that after a long time debt will converge to zero. Nevertheless it is still interesting to look at the internal dynamics of the model and whether they are similar to what is observed in the data.}. \\

In table \ref{tab:std} we can see the standard deviation of the variables real quarterly gross domestic product, real quarterly fixed capital formation, real quarterly household consumption, and quarterly employment. The variables are first filtered with the Hodrick-Prescott filter and then normalized with their own trend such that the variables can be compared at the same scale. We can see that all the models generate too little variation for GDP, investment, and inflation but are quite close. All models do get the ordering of variance correctly with investment being the most volatile, then in order, GDP, Consumption and finally inflation. We can see that EH-ABM and OC-ABM improve on the baseline ABM by decreasing consumption volatility in comparison to GDP volatility.
\mycomment{
In figure \ref{fig:cycles-abm}
we can see the generated paths of the agent based model calibrated in the first quarter of 2012. We can see that the variation of household consumption is quite close to the variation of GDP whereas the variation of real capital formation is larger than that of GDP. \\
}

\begin{table}[h!]
    \centering
    \caption{Comparison of variance of variables}
    \begin{tabular}{c|cccc}
    \toprule
         &GDP &Investment&Consumption& Inflation \\
         \midrule 
        Data &0.0122&0.0528& 0.0086 & 0.0082 \\ \midrule
        ABM &0.0098& 0.0371 &0.0091&0.0066\\
        EH-ABM &0.0098&0.0382&0.0089&0.0067\\
        OC-ABM &0.00900 &0.0388 & 0.0076& 0.0066\\
    \bottomrule
    \end{tabular}
    	\begin{tablenotes}
		\footnotesize \item \textit{Note}: Standard deviation of the variables real quarterly gross domestic product, real quarterly fixed capital formation, real quarterly household consumption, and quarterly employment. 
	\end{tablenotes}
    \label{tab:std}
\end{table}

\begin{figure}[!h]
    \centering
    \includesvg[width=\textwidth]{figs/optimal_consumption/calibrated/cycles_abm_uf.svg}
    \caption{Normalized business cycles of OC-ABM generated paths of variables real quarterly gross domestic product, GDP deflator growth, real quarterly household consumption, and real quarterly capital formation. Generated variable paths are first filtered with the Hodrick-Prescott filter and normalized by dividing by the trend path. Monte Carlo sample paths are generated for 2012Q1 for 25 times.}
    \label{fig:cycles-abm}
\end{figure}
Figure \ref{fig:autocorr-abm} shows the auto-correlations\footnote{We refer to the sample auto-correlation function which is defined as 
$r_k=\frac{\frac{1}{T} \sum_{t=1}^{T-k}\left(y_t-\bar{y}\right)\left(y_{t+k}-\bar{y}\right)}{\frac{1}{T} \sum_{t=1}^{T}\left(y_t-\bar{y}\right)\left(y_{t}-\bar{y}\right)}
$.} of the cyclical components of real gross domestic product at time $t$ and the cyclical components of a number of variables. The auto-correlations are calculated by using the generated Monte Carlo for 2010Q1:2012Q4. We can observe that for all variables the auto-correlations of the variables look similar to the auto-correlation of the data\footnote{The auto- and cross correlation calculated from the data are also estimates meaning t}. We note that there is slower decay in the auto-correlation of wages and household consumption as compared to the data. This suggests that there are not enough frictions in the job-searching mechanism. It is interesting for future work to add these frictions. \\ 

Figure \ref{fig:crosscorr-abm} shows the cross-correlations\footnote{We refer to the sample cross-correlation function which is defined as 
\\$\widehat{R}_{x y}(m)= \begin{cases}\sum_{n=0}^{N-m-1} x_{n+m} y_n^*, & m \geq 0, \\ \widehat{R}_{y x}^*(-m), & m<0 .\end{cases}$. We normalize the cross-correlation function numerically.} between the cyclical components of real gross domestic product at time $t$ and the cyclical components of a number of variables at time $t - lag$. Firstly, the cross-correlations of real capital formation and of operating surplus look quite similar to the cross-correlations observed in the data. We also observe that the cross-correlations of wages is pro-cyclical in the data whereas in the sample paths generated by the ABM wages is coincident with GDP. The reason that wages are coincidental with GDP is most likely due to the mechanism that wages increase when production also increases due to overtime payment. In future work, we could change the job matching mechanism such that filling a vacancy takes time such as for instance a quarter instead of it being instantaneously such as the model is currently specified. The GDP deflator cross-correlation match the least up with the data generated cross-correlations but the signs are correct. It could be that the price setting mechanism is too simplistic or that the expectation formation model is too simplistic. Note however that the empirical cross- and auto-correlations are also estimates and their actual values might be closer or further away than our modelled counterparts.\\ 

In figure \ref{fig:crosscorr-oc-abmx-uf}, we can see the cross correlations of the generated variable paths of unconditional ARX optimal consumption agent-based model. We can see that the cross-correlations of the GDP deflator and real GDP of the model are much closer to the data estimated cross-correlations. This indicates that the ARX agent based model generates more dynamics that are more internally consistent than the AR agent based model. This suggest that ARX expectations is a more realistic assumption than AR expectations.

\begin{figure}
    \centering
    \includesvg[width=\textwidth]{figs/optimal_consumption/calibrated/autocorr_abm_uf.svg}
		\caption{Auto-correlations of quarterly real gross domestic product of of paths of variables the real quarterly gross domestic product, real quarterly fixed capital formation, real quarterly household consumption, and quarterly employment. Generated variable paths are first filtered with the Hodrick-Prescott filter. Monte Carlo sample paths are generated for 2012Q1:2016Q4 for 25 times. The blue line is the mean of the calculated auto-correlation whereas the blue error-bar is the standard deviation of the calculated paths. The red line is the auto-correlation observed in the data. }
    \label{fig:autocorr-abm}
\end{figure}


\begin{figure}[!h]
    \centering
      \includesvg[width=\textwidth]{figs/optimal_consumption/calibrated/crosscorr_abm_uf.svg}
    \caption{Cross-correlations of quarterly real gross domestic product of the generated paths of the variables real quarterly fixed capital formation, real quarterly household consumption, and quarterly employment. Generated variable paths are first filtered with the Hodrick-Prescott filter. Monte Carlo sample paths are generated for 2012Q1:2016Q4 for 25 times. The blue line is the mean of the calculated cross-correlation whereas the blue error-bar is the standard deviation of the calculated paths. The red line is the cross-correlation observed in the data.}
    \label{fig:crosscorr-abm}
\end{figure}

